\chapter{Methods and Software Implementation}
\section{Entry-aware splitting}
For a candidate split at a node, let $t_1,\ldots,t_K$ be its distinct event
times. Let $d_k$ and $d_{1k}$ denote event counts in the parent and left
daughter, and let $Y_k$ and $Y_{1k}$ denote their risk-set sizes. The
delayed-entry log-rank numerator and variance term are
\begin{align}
 U&=\sum_{k:Y_k>0}\left(d_{1k}-d_k\frac{Y_{1k}}{Y_k}\right),\\
 V&=\sum_{k:Y_k>1}
 \frac{Y_{1k}(Y_k-Y_{1k})d_k(Y_k-d_k)}{Y_k^2(Y_k-1)}.
\end{align}
Candidate splits are compared using the magnitude of the standardized
statistic when its variance is positive. Entry awareness changes the
risk-set counts through Equation~\eqref{eq:risk}; observed event counts are
unchanged. These are the risk-set modifications specified in the
professor-provided project outline \citep{projectoutline}.

\section{Terminal-node estimates and returned predictions}
For terminal node $h$, let $d_h(t_k)$ be its event count and $Y_h(t_k)$ its
entry-aware risk-set size. The node estimates are
\begin{align}
 \widehat H_h(t)&=\sum_{t_k\leq t}\frac{d_h(t_k)}{Y_h(t_k)},\\
 \widehat S_h(t)&=\prod_{t_k\leq t}
       \left(1-\frac{d_h(t_k)}{Y_h(t_k)}\right).
\end{align}
Only node event times with a nonzero risk set contribute. A new subject is
routed to node $h_b(x)$ in tree $b$ using its covariates. The package returns
averages of the tree-specific cumulative hazards and product-limit curves:
\begin{align}
 \widehat H_e(t\mid x)&=\frac{1}{B}\sum_{b=1}^{B}\widehat H_{h_b(x)}(t),\\
 \widehat S_e(t\mid x)&=\frac{1}{B}\sum_{b=1}^{B}\widehat S_{h_b(x)}(t).
\end{align}
OOB predictions average over the trees that excluded the training subject.

The outline also describes $\exp\{-\widehat H_e(t\mid x)\}$ as a possible
survival construction. This is not generally identical to the average
product-limit curve returned by the package. The experiments retain the
package's returned survival curves; they do not silently replace them by
the exponential of the averaged cumulative hazard. This construction should
be confirmed when finalizing the thesis specification.

A test subject's entry time is not needed to route it through the trees.
Entry becomes relevant when conditioning a prediction on observation or
evaluating it against an outcome.

\section{Integration, validation, and supported scope}
The user supplies one entry time per original input row through
\texttt{rfsrc(..., entry.time = L)}. The interface aligns entries with rows
retained after missing-data processing and checks that retained entries are
finite, nonnegative, and no later than follow-up. Positive entry with a
missing follow-up time is rejected; missing predictors remain imputable.
The C implementation uses a lower-bound lookup to map entries onto the
master time grid, without requiring entries to coincide with exit times.

Entry times are stored on the grow object and saved forest, and are passed
explicitly during restoration and prediction. Refitting on subsets or
bootstrap samples subsets entries with the observations. Tuning aligns
training and holdout entries and uses holdout subjects' entries for the LT
concordance comparison. Training outcomes used to estimate nuisance
distributions are kept separate from evaluation outcomes.

Positive entries are supported with the single-event default log-rank rule
and with outcome-independent random splitting. Other rules, including
log-rank-score, Brier-gradient, and custom survival rules, are rejected with
positive entries because their split statistics have not been fully adapted.
Allowing random splits does not establish general LT validity for every
outcome family. The thesis evaluates single-event log-rank forests.

For nonnegative follow-up, setting all entries to zero reduces
Equation~\eqref{eq:risk} to the ordinary risk indicator. Ordinary calls can
then use the existing fast concordance algorithms. Positive-entry
concordance currently uses pairwise enumeration. Zero-entry agreement is a
necessary compatibility check, not a proof of correctness for positive entry.

\section{Evaluation targets}
Truth-based curve error evaluates $\widehat S(t\mid X)$ against the
source-population survival function. Outcome-based evaluation instead uses
observed training subjects. In the experiments, RSF uses OOB curves and Cox
uses in-sample fitted curves. Therefore, outcome-score differences between
the two model families are not comparisons under identical validation designs.

\subsection{Concordance}
For an observed event at $U_i<U_j$, the implemented LT comparison requires
$L_j\leq U_i$. The simulation risk marker is
\[
 r_i=\sum_{g=1}^{G}\{1-\widehat S(t_g\mid X_i)\}.
\]
The reported concordance error is $1-C$ using the package's pair and
tie conventions, with the delayed-entry restriction. The simulation calls
the unweighted helper. The package's native performance machinery also
supports its existing Uno-weighted option; that is not the weighting used
for the simulation C-index tables.

\subsection{Observation-conditioned Brier score}
Section 9.2 of the outline \citep{projectoutline} requires censoring adjustment, restriction to
$L_i\leq t$, and conditioning on observation. It does not supply an expanded
estimator. The following construction is our interpretation of that
instruction and remains subject to the professor's scientific review.

Related methodological references consulted during development were the
\texttt{JMbayes2} documentation for dynamic Brier scoring
\citep{jmbayes2brier} and the supplement to the Fu--Simonoff survival-tree
study \citep[Section~4.1]{fuSimonoffSupplement}.
The former documents prediction from a specified starting time to a later
horizon, with IPCW and left-truncation support. Our score instead averages
subject-specific entry-conditioned predictions over subjects entered by $t$.
The latter presents the ordinary censoring-adjusted Brier formula, but its
reported test-set evaluation uses uncensored event times.
These references provide context; neither supplies the exact estimator below.

With the implemented inclusive-entry convention, the prediction after
entry is
\[
 q_i(t)=\frac{\widehat S(t\mid X_i)}
                  {\widehat S(L_i^-\mid X_i)},\qquad t\geq L_i,
\]
when the denominator is positive. Let $G_i$ denote censoring survival under
the chosen censoring model. Let $\mathcal E_t$ contain the selected subjects with $L_i\leq t$ and a
finite conditional prediction, and write $n_t=|\mathcal E_t|$. The score is
\begin{align}
 \widehat{\mathrm{BS}}_{\mathrm{entry}}(t)
 =\frac{1}{n_t}\sum_{i\in\mathcal E_t}
 \Bigg[&
 \delta_i\ind(U_i\leq t)q_i(t)^2
        \frac{\widehat G_i(L_i^-)}{\widehat G_i(U_i^-)}
 \nonumber\\
 &+\ind(U_i>t)\{1-q_i(t)\}^2
        \frac{\widehat G_i(L_i^-)}{\widehat G_i(t)}
 \Bigg].
\end{align}
Here $n_t$ counts those entered subjects with finite predictions. Censoring
by $t$ gives a zero contribution while retaining the subject in this
denominator. Subjects not yet entered are excluded. No empirical entry-CDF
weight is used: the entry distribution among retained subjects generally
differs from its source-population distribution.

The experiments use a pooled product-limit censoring estimate with
delayed-entry risk sets. Failures are removed before censoring at tied exit
times. Valid IPCW interpretation requires independent censoring given the
information in the censoring model and positive observation probabilities.
A pooled model additionally assumes a common censoring hazard. Covariate-
or entry-dependent censoring in some scenarios does not automatically
satisfy that stronger assumption; those scores are descriptive under this
working censoring model.

The scalar reported as Brier/CRPS is the trapezoidal approximation to
\[
 \mathrm{CRPS}=\int_{t_{\min}}^{t_{\max}}
                  \widehat{\mathrm{BS}}_{\mathrm{entry}}(t)\,dt.
\]
It is an unnormalized time integral, not an average bounded by one.
Horizon lengths vary across scenarios. A grid with no entrants yields an
undefined score; it is not filled with zero.

\subsection{Additional time-dependent AUC metric}
The professor's outline does not prescribe an AUC estimator. Our added
cumulative/dynamic AUC uses entered observed failures as cases and entered
subjects with $U_j>t$ as controls. Its marker is $M_i(t)=1-q_i(t)$, with
weights
\[
 w_i^{\mathrm{case}}(t)=
 \frac{\widehat G_i(L_i^-)}{\widehat G_i(U_i^-)},\qquad
 w_j^{\mathrm{control}}(t)=
 \frac{\widehat G_j(L_j^-)}{\widehat G_j(t)}.
\]
AUC is the weighted case--control fraction for which the case has a larger
marker, with half credit for ties. The reported error is
$1-\overline{\mathrm{AUC}}$, where the bar denotes trapezoidal time
averaging over defined grid points. This is a descriptive time average,
not an event-distribution-weighted population integrated AUC, and the
current averaging can bridge undefined interior points. The censoring
assumptions are the same as for Brier scoring.

These entered-cohort scores do not estimate time-zero prediction error
for unobserved source-population subjects. Their targets differ from
truth-based MAD. External scoring requires evaluation outcomes and entries;
unlabeled predictions cannot be assigned an outcome-based score. The
native new-data prediction interface has no test-entry channel, so its
built-in error output should not be substituted for explicit LT test-set
scoring.
